\documentclass[11pt]{article}

\usepackage[utf8]{inputenc}
\usepackage[T1]{fontenc}
\usepackage{lmodern}
\usepackage{geometry}
\usepackage{amsmath,amssymb,amsfonts,amsthm,mathtools}
\usepackage{bm}
\usepackage{enumitem}
\usepackage{microtype}
\usepackage{hyperref}
\usepackage{titlesec}
\usepackage{booktabs}
\usepackage{array}
\usepackage{setspace}
\usepackage{mathrsfs}
\usepackage{physics}

\geometry{margin=1in}
\hypersetup{colorlinks=true, linkcolor=black, urlcolor=blue, citecolor=black}
\setlist[enumerate]{topsep=0.4em,itemsep=0.35em}
\setlist[itemize]{topsep=0.3em,itemsep=0.25em}
\titleformat{\section}{\large\bfseries}{\thesection.}{0.5em}{}
\titleformat{\subsection}{\normalsize\bfseries}{\thesubsection.}{0.5em}{}
\setstretch{1.06}

\newcommand{\Aether}{\AE{}ther}
\newcommand{\Aflow}{\AE{}ther-flow}
\newcommand{\TheoryName}{\Aflow{} Interpretation of Relativity}
\newcommand{\TheoryProgram}{\Aether{} / \Aflow{} framework}

\newcommand{\CompletedMetric}{g^{\sharp}}
\newcommand{\MatterSpace}{\Psi}

\newcommand{\Car}{\mathrm{car}}
\newcommand{\Cur}{\mathrm{cur}}
\newcommand{\Hom}{\mathrm{hom}}

\newcommand{\ForSpace}[1]{\mathcal I_{#1}^{\mathrm{for}}}
\newcommand{\PrimShell}{\mathcal S_{\mathrm{prim}}}
\newcommand{\MetricOut}{\mathscr G_{\mathrm{out}}}
\newcommand{\MatterOut}{\mathscr T_{\mathrm{out}}}

\newcommand{\PreCoreTransportCore}{\mathfrak S^{\mathrm{tr}}}
\newcommand{\PreCoreTrBody}[1]{\mathcal B_{#1}^{\mathrm{precore,tr}}}
\newcommand{\PreCoreTrShell}[1]{\mathcal Z_{#1}^{\mathrm{precore,tr}}}

\newcommand{\SubPreSectionCore}{\mathfrak U^{\mathrm{sec}}}
\newcommand{\SubSectionPackage}{\mathfrak V^{\mathrm{subsec}}}
\newcommand{\SubSectionSpace}[1]{\mathcal M_{#1}^{\mathrm{subsec}}}
\newcommand{\SubSectionBody}[1]{\mathcal B_{#1}^{\mathrm{subsec}}}
\newcommand{\SubSectionProj}[1]{\Pi_{#1}^{\mathrm{subsec}}}
\newcommand{\SubSectionFiber}[2]{\mathcal F_{#1}^{\mathrm{subsec}}(#2)}
\newcommand{\SubSectionFunctional}[1]{\mathscr E_{#1}^{\mathrm{subsec}}}
\newcommand{\SubSectionShell}[1]{\mathcal Z_{#1}^{\mathrm{subsec}}}

\newcommand{\MicroSectionPackage}{\mathfrak W^{\mathrm{microsec}}}
\newcommand{\MicroSectionSpace}[1]{\mathcal N_{#1}^{\mathrm{microsec}}}
\newcommand{\MicroSectionBody}[1]{\mathcal B_{#1}^{\mathrm{microsec}}}
\newcommand{\MicroSectionProj}[1]{\Pi_{#1}^{\mathrm{microsec}}}
\newcommand{\MicroSectionFiber}[2]{\mathcal F_{#1}^{\mathrm{microsec}}(#2)}
\newcommand{\MicroSectionSubMin}[1]{\mathfrak m_{#1}^{\mathrm{microsec}}}
\newcommand{\MicroSectionFunctional}[1]{\mathscr E_{#1}^{\mathrm{microsec}}}
\newcommand{\MicroSectionShell}[1]{\mathcal Z_{#1}^{\mathrm{microsec}}}

\newtheorem{definition}{Definition}
\newtheorem{proposition}{Proposition}
\newtheorem{remark}{Remark}
\newtheorem{corollary}{Corollary}
\newtheorem{theorem}{Theorem}

\title{\textbf{The \TheoryName{}}\\[0.4em]
\large Primitive-Reservoir Observer-Reduction\\
\large Exact-Shell-Transport-Section-Generating Substrate Sub-Section Dynamics\\
\large from Sub-Section-Generating Substrate Micro-Section Dynamics}
\author{}
\date{}

\begin{document}

\maketitle

\begin{abstract}
The current primitive-reservoir observer-reduction frontier already moved the live burden below the observer-relevant sub-pre-core exact-shell transport-section core itself. The section-core theorem showed that the current core is canonically generated by a deeper exact-shell-transport-section-generating substrate sub-section dynamics package \(\SubSectionPackage\). The accompanying routing note then fixed the next honest move: either derive or justify that current sub-section package from still more primitive substrate structure, or prove a genuinely smaller theorem-level collapse strictly inside it.

The present note takes the preferred first route. For each sector it introduces \emph{sub-section-generating substrate micro-section dynamics}: a compact convex micro-section body over the recorded sub-section body, an affine projection to that fixed body, a nonnegative smooth micro-section functional with unique fiberwise minimizers, and an exact micro-section shell projecting diffeomorphically to the recorded exact sub-section shell. The current sub-section functional is then no longer an unexplained primitive datum. It is the effective fiberwise minimum value induced by that deeper micro-section law, and the current sub-section shell is the projected exact shell of the deeper package.

The benchmark boundary remains fixed throughout. The one operative metric is still exactly \(\CompletedMetric\), the ordinary matter route is unchanged, there is no second operative metric, the low-energy matter law is not enlarged, and no stronger promotion language is licensed. The positive result is therefore narrow but benchmark facing: the exact-shell-transport-section-generating substrate sub-section dynamics are no longer left as the primitive stopping point on the active line. The remaining honest burden moves one level deeper again, to the sub-section-generating substrate micro-section dynamics themselves or to a sharper collapse strictly inside that deeper package.
\end{abstract}

\section{Introduction}

The active derivational program of \TheoryName{} has now reached a sharper burden than another same-output relay. The current active-primary theorem already showed that the observer-relevant sub-pre-core exact-shell transport-section core \(\SubPreSectionCore\) is not primitive: it is the canonical section selected by a deeper exact-shell-transport-section-generating substrate sub-section dynamics package \(\SubSectionPackage\). The accompanying routing note then fixed the consequence of that theorem. The next honest benchmark-facing manuscript must either derive or justify \(\SubSectionPackage\) itself from still more primitive substrate structure, or prove a genuinely smaller theorem-level collapse strictly inside that package.

The present note takes the direct derivation route. Its claim is not that final substrate microphysics has been identified, and it is not that the benchmark exact-relativistic package has changed. Its claim is narrower and benchmark facing. The current sub-section package is shown to arise canonically from a more primitive micro-section layer whose projected exact-shell transport already contains the benchmark-relevant structure of \(\SubSectionPackage\). In that precise sense the current sub-section dynamics are no longer the primitive place where the active line has to stop.

Two features matter here. First, the theorem targets the current sub-section package itself rather than re-spending the already settled section-core handoff, so it does not count progress merely by reproducing the same generated section core from yet another relay. Second, the deeper micro-section package is not merely a renaming of the current sub-section package. It adds an explicit micro-section state space, a micro-section body, a micro-to-sub-section projection, and a deeper fiberwise minimizing law whose effective energy reproduces the recorded sub-section functional on the benchmark-relevant body.

\paragraph{Framework and claim boundary.}
\TheoryProgram{} denotes the broader \Aether{} / \Aflow{} framework built on the claims that the \Aether{} is the underlying four-dimensional substrate of reality, the \Aflow{} is its intrinsic ordered motion, observed three-dimensional space is the local experiential slice of that deeper substrate, S-time is the experienced order of change arising from matter, light, and the \Aflow{}, and observed expansion is the three-dimensional appearance of deeper four-dimensional ordered motion. Within that framework, \TheoryName{} denotes the exact-closure theory in which the effective gravitational dynamics are adopted to be exactly Einsteinian with universal matter coupling. In the active sequence, \emph{adoption} means use of that established relativistic dynamics without claiming substrate derivation, while \emph{derivation} is reserved for a first-principles recovery from explicit substrate variables.

The present note stays strictly inside that boundary. It does not alter the benchmark exact-closure package. It does not introduce a second operative metric or a new low-energy matter law. It only proves that the current exact-shell-transport-section-generating substrate sub-section dynamics can itself be generated from a more primitive micro-section-level substrate dynamics.

\section{Fixed Inputs}

\begin{definition}[Recorded primitive continuation data]
\label{def:fixed}
Fix the following continuation data from the active primitive-reservoir observer-reduction line.
\begin{enumerate}
    \item For each sector label \(\sigma\in\{\Car,\Cur,\Hom\}\), a primitive forcing shell
    \[
    \ForSpace{\sigma}.
    \]
    \item The product shell
    \[
    \PrimShell
    \equiv
    \ForSpace{\Car}\times \ForSpace{\Cur}\times \ForSpace{\Hom}.
    \]
    \item The recorded metric and ordinary-matter outputs
    \[
    \MetricOut:\PrimShell\to\{\text{observer metrics}\},
    \qquad
    \MatterOut:\PrimShell\times\MatterSpace\to\{\text{observer stress tensors}\},
    \]
    satisfying
    \begin{equation}
    \MetricOut(\mathbf w)=\CompletedMetric,
    \qquad
    \MatterOut(\mathbf w,\psi)
    =
    T_{\mu\nu}^{\mathrm{matter}}[\CompletedMetric,\psi]
    \label{eq:fixedoutputs}
    \end{equation}
    for every \(\mathbf w\in\PrimShell\) and every matter configuration \(\psi\in\MatterSpace\).
\end{enumerate}
\end{definition}

\begin{definition}[Recorded transport-core data]
\label{def:transport}
Fix \(\sigma\in\{\Car,\Cur,\Hom\}\). The active line already fixes the benchmark-relevant transport-core data:
\begin{enumerate}
    \item a finite-dimensional pre-core exact-shell transport body
    \[
    \PreCoreTrBody{\sigma};
    \]
    \item a closed embedded exact-shell transport shell
    \[
    \PreCoreTrShell{\sigma}\subset\PreCoreTrBody{\sigma};
    \]
    \item the already recorded downstream shell transport from \(\PreCoreTrShell{\sigma}\) to the observer-relevant pre-nucleus exact-shell core, the current transport nucleus, and the fixed proto-germ exact-shell target.
\end{enumerate}
\end{definition}

\begin{remark}
Definitions \ref{def:fixed} and \ref{def:transport} are fixed inputs. The one operative metric, the ordinary matter route, the transport-core body \(\PreCoreTrBody{\sigma}\), the exact transport shell \(\PreCoreTrShell{\sigma}\), and the already recorded downstream shell transport are not reopened here. The present note asks only whether the current sub-section dynamics can themselves be generated from still more primitive substrate structure.
\end{remark}

\section{Recorded Sub-Section Target}

\begin{definition}[Recorded exact-shell-transport-section-generating substrate sub-section dynamics]
\label{def:subsec}
Fix \(\sigma\in\{\Car,\Cur,\Hom\}\). The recorded exact-shell-transport-section-generating substrate sub-section dynamics package \(\SubSectionPackage\) for sector \(\sigma\) consists of the following data.
\begin{enumerate}
    \item A finite-dimensional Euclidean sub-section state space \(\SubSectionSpace{\sigma}\), a compact convex sub-section body
    \[
    \SubSectionBody{\sigma}\subset\SubSectionSpace{\sigma}
    \]
    with nonempty interior, and an affine surjection
    \[
    \SubSectionProj{\sigma}:\SubSectionSpace{\sigma}\to\PreCoreTrBody{\sigma}.
    \]
    For each \(w\in\PreCoreTrBody{\sigma}\), write
    \[
    \SubSectionFiber{\sigma}{w}
    \equiv
    \bigl(\SubSectionProj{\sigma}\bigr)^{-1}(w)\cap\SubSectionBody{\sigma}.
    \]
    Each such fiber is required to be nonempty, compact, and convex.
    \item A nonnegative smooth fiber-selection functional
    \[
    \SubSectionFunctional{\sigma}:\SubSectionSpace{\sigma}\to[0,\infty)
    \]
    such that, for every \(w\in\PreCoreTrBody{\sigma}\), the restriction of \(\SubSectionFunctional{\sigma}\) to \(\SubSectionFiber{\sigma}{w}\) is strictly convex and attains a unique minimizer. The resulting minimizer depends smoothly on \(w\).
    \item A closed embedded exact sub-section shell
    \[
    \SubSectionShell{\sigma}\subset\SubSectionBody{\sigma}
    \]
    satisfying
    \[
    \SubSectionShell{\sigma}
    =
    \left\{
    y\in\SubSectionBody{\sigma}:
    \SubSectionProj{\sigma}(y)\in\PreCoreTrShell{\sigma},
    \;
    \SubSectionFunctional{\sigma}(y)
    =
    \min_{z\in\SubSectionFiber{\sigma}{\SubSectionProj{\sigma}(y)}}
    \SubSectionFunctional{\sigma}(z)
    \right\},
    \]
    and
    \[
    \SubSectionProj{\sigma}\big|_{\SubSectionShell{\sigma}}
    :
    \SubSectionShell{\sigma}
    \xrightarrow{\cong}
    \PreCoreTrShell{\sigma}
    \]
    is a diffeomorphism.
    \item Benchmark faithfulness, meaning preservation of the benchmark outputs \eqref{eq:fixedoutputs}, no second operative metric, no enlarged low-energy matter law, and no premature promotion from adoption to completed derivation.
\end{enumerate}
\end{definition}

\begin{remark}
Definition \ref{def:subsec} is the fixed target already carried by the current active-primary theorem. The present note does not spend its move on another theorem about the already generated section core. It records the current sub-section package only because the deeper micro-section dynamics are defined relative to it.
\end{remark}

\section{Sub-Section-Generating Substrate Micro-Section Dynamics}

\begin{definition}[Sub-section-generating substrate micro-section dynamics]
\label{def:microsec}
Fix \(\sigma\in\{\Car,\Cur,\Hom\}\). A sub-section-generating substrate micro-section dynamics package \(\MicroSectionPackage\) for sector \(\sigma\) consists of the following data.
\begin{enumerate}
    \item A finite-dimensional Euclidean micro-section state space \(\MicroSectionSpace{\sigma}\), a compact convex micro-section body
    \[
    \MicroSectionBody{\sigma}\subset\MicroSectionSpace{\sigma}
    \]
    with nonempty interior, and an affine surjection
    \[
    \MicroSectionProj{\sigma}:\MicroSectionSpace{\sigma}\to\SubSectionSpace{\sigma}
    \]
    satisfying
    \[
    \MicroSectionProj{\sigma}\bigl(\MicroSectionBody{\sigma}\bigr)=\SubSectionBody{\sigma}.
    \]
    For each \(y\in\SubSectionBody{\sigma}\), write
    \[
    \MicroSectionFiber{\sigma}{y}
    \equiv
    \bigl(\MicroSectionProj{\sigma}\bigr)^{-1}(y)\cap\MicroSectionBody{\sigma}.
    \]
    Each such fiber is required to be nonempty, compact, and convex.
    \item A nonnegative smooth micro-section functional
    \[
    \MicroSectionFunctional{\sigma}:\MicroSectionSpace{\sigma}\to[0,\infty)
    \]
    such that, for every \(y\in\SubSectionBody{\sigma}\), the restriction of \(\MicroSectionFunctional{\sigma}\) to \(\MicroSectionFiber{\sigma}{y}\) is strictly convex and attains a unique minimizer. The resulting minimizer depends smoothly on \(y\); write it as
    \[
    \MicroSectionSubMin{\sigma}:\SubSectionBody{\sigma}\to\MicroSectionBody{\sigma}.
    \]
    These data satisfy
    \[
    \MicroSectionProj{\sigma}\circ\MicroSectionSubMin{\sigma}
    =
    \mathrm{id}_{\SubSectionBody{\sigma}},
    \]
    and the effective minimum value recovers the recorded sub-section functional on the benchmark-relevant body:
    \[
    \SubSectionFunctional{\sigma}(y)
    =
    \MicroSectionFunctional{\sigma}\bigl(\MicroSectionSubMin{\sigma}(y)\bigr)
    =
    \min_{u\in\MicroSectionFiber{\sigma}{y}}
    \MicroSectionFunctional{\sigma}(u)
    \qquad
    \text{for every }
    y\in\SubSectionBody{\sigma}.
    \]
    \item A closed embedded exact micro-section shell
    \[
    \MicroSectionShell{\sigma}\subset\MicroSectionBody{\sigma}
    \]
    satisfying
    \[
    \MicroSectionShell{\sigma}
    =
    \left\{
    u\in\MicroSectionBody{\sigma}:
    \begin{array}{l}
    \MicroSectionProj{\sigma}(u)\in\SubSectionShell{\sigma},\\[0.2em]
    \MicroSectionFunctional{\sigma}(u)
    =
    \displaystyle\min_{v\in\MicroSectionFiber{\sigma}{\MicroSectionProj{\sigma}(u)}}
    \MicroSectionFunctional{\sigma}(v)
    \end{array}
    \right\},
    \]
    and
    \[
    \MicroSectionProj{\sigma}\big|_{\MicroSectionShell{\sigma}}
    :
    \MicroSectionShell{\sigma}
    \xrightarrow{\cong}
    \SubSectionShell{\sigma}
    \]
    is a diffeomorphism.
    \item Benchmark faithfulness, meaning preservation of the benchmark outputs \eqref{eq:fixedoutputs}, no second operative metric, no enlarged low-energy matter law, and presentation as a deeper substrate-side source for the current sub-section package rather than as a replacement benchmark theory.
\end{enumerate}
\end{definition}

\begin{remark}
The label \emph{micro-section} is used here only as a local marker for a deeper layer below the current sub-section frontier. It does not by itself assert a completed microscopic ontology or a final primitive law. Its role is only to name the next sharper package below the current sub-section dynamics.
\end{remark}

\section{Induced Recovery of the Current Sub-Section Package}

\begin{proposition}[The effective sub-section functional is induced by fiberwise micro minimization]
\label{prop:functional}
Fix a benchmark-faithful micro-section package. Then the effective minimum-value functional
\[
y\longmapsto
\min_{u\in\MicroSectionFiber{\sigma}{y}}
\MicroSectionFunctional{\sigma}(u)
\]
is smooth on \(\SubSectionBody{\sigma}\) and agrees there with the recorded sub-section functional \(\SubSectionFunctional{\sigma}\).
\end{proposition}

\begin{proof}
Definition \ref{def:microsec}(2) states that the restriction of \(\MicroSectionFunctional{\sigma}\) to each fiber \(\MicroSectionFiber{\sigma}{y}\) has a unique minimizer and that this minimizer depends smoothly on \(y\). Hence the minimum-value map is smooth and equals
\[
\MicroSectionFunctional{\sigma}\bigl(\MicroSectionSubMin{\sigma}(y)\bigr).
\]
The same definition requires this quantity to equal the recorded \(\SubSectionFunctional{\sigma}(y)\) for every \(y\in\SubSectionBody{\sigma}\).
\end{proof}

\begin{proposition}[Projected exact-shell transport recovers the recorded sub-section shell]
\label{prop:shell}
Fix a benchmark-faithful micro-section package. Then
\[
\MicroSectionProj{\sigma}\bigl(\MicroSectionShell{\sigma}\bigr)
=
\SubSectionShell{\sigma},
\]
and
\[
\MicroSectionProj{\sigma}\big|_{\MicroSectionShell{\sigma}}
:
\MicroSectionShell{\sigma}\xrightarrow{\cong}\SubSectionShell{\sigma}
\]
is a diffeomorphism.
\end{proposition}

\begin{proof}
The displayed diffeomorphism is part of Definition \ref{def:microsec}(3). Its image is therefore exactly \(\SubSectionShell{\sigma}\).
\end{proof}

\begin{proposition}[Micro-section dynamics canonically reproduce the current sub-section package]
\label{prop:package}
Fix a benchmark-faithful micro-section package. Then the current exact-shell-transport-section-generating substrate sub-section dynamics of Definition \ref{def:subsec} are exactly the effective benchmark-relevant package induced on the recorded body \(\SubSectionBody{\sigma}\) by the deeper micro-section dynamics. In particular:
\begin{enumerate}
    \item the recorded sub-section body and projection remain unchanged;
    \item the recorded sub-section functional on \(\SubSectionBody{\sigma}\) is the fiberwise minimum-value functional induced by the deeper micro-section package;
    \item the recorded exact sub-section shell is the projected exact shell of the deeper micro-section package.
\end{enumerate}
\end{proposition}

\begin{proof}
Item (1) is part of the fixed target data in Definition \ref{def:subsec} together with Definition \ref{def:microsec}(1). Item (2) is Proposition \ref{prop:functional}. Item (3) is Proposition \ref{prop:shell}. These statements are exactly the benchmark-relevant constituents of the current sub-section package.
\end{proof}

\begin{proposition}[The benchmark package remains fixed]
\label{prop:boundary}
For every benchmark-faithful micro-section package, the operative metric remains exactly \(\CompletedMetric\), the ordinary matter route remains unchanged, no second operative metric is introduced, and the low-energy matter law is not enlarged.
\end{proposition}

\begin{proof}
This is part of the benchmark-faithfulness requirement in Definition \ref{def:microsec}(4). Equation \eqref{eq:fixedoutputs} remains unchanged on the full primitive forcing shell, so the deeper micro-section package does not alter the benchmark exact-closure package.
\end{proof}

\begin{corollary}[The current section-core theorem applies unchanged]
\label{cor:section}
Fix a benchmark-faithful micro-section package. Because Proposition \ref{prop:package} reproduces exactly the recorded sub-section package of Definition \ref{def:subsec}, the previously recorded theorem deriving the observer-relevant sub-pre-core exact-shell transport-section core \(\SubPreSectionCore\) from exact-shell-transport-section-generating substrate sub-section dynamics applies unchanged. Consequently the current section core and all previously recorded downstream shell-transport consequences remain available unchanged.
\end{corollary}

\begin{proof}
The cited theorem requires exactly the sub-section package of Definition \ref{def:subsec} as its live input at current scope. Proposition \ref{prop:package} reproduces that package on the nose, so the same section-core recovery and downstream consequences follow without modification.
\end{proof}

\section{Main Theorem}

\begin{theorem}[Exact-shell-transport-section-generating substrate sub-section dynamics from sub-section-generating substrate micro-section dynamics]
\label{thm:main}
Fix the primitive continuation data of Definition \ref{def:fixed}, the recorded transport-core data of Definition \ref{def:transport}, and the recorded exact-shell-transport-section-generating substrate sub-section dynamics of Definition \ref{def:subsec}. Then, for each sector \(\sigma\in\{\Car,\Cur,\Hom\}\), every benchmark-faithful sub-section-generating substrate micro-section dynamics package canonically produces the current sub-section package as the effective fiberwise minimum-value output of the deeper micro-section layer. More precisely:
\begin{enumerate}
    \item the current exact-shell-transport-section-generating substrate sub-section dynamics are reproduced unchanged at benchmark scope by the deeper micro-section dynamics;
    \item because the induced sub-section package is exactly the one already used by the current section-core theorem, the observer-relevant sub-pre-core exact-shell transport-section core \(\SubPreSectionCore\) and the previously recorded downstream consequences remain available unchanged;
    \item the same benchmark one-metric package remains fixed: the operative metric is still exactly \(\CompletedMetric\), the ordinary matter route is unchanged, there is no second operative metric, and the low-energy matter law is not enlarged;
    \item consequently the remaining honest burden on the active line is deeper again: derive or justify the sub-section-generating substrate micro-section dynamics themselves from still more primitive substrate structure, or else prove a genuinely smaller theorem-level collapse strictly inside that deeper package.
\end{enumerate}
\end{theorem}

\begin{proof}
Item (1) is Proposition \ref{prop:package}. Item (2) is Corollary \ref{cor:section}. Item (3) is Proposition \ref{prop:boundary}. Item (4) is the project-level consequence of the previous items together with the routing safeguard against counting another same-output deeper-origin relay as new front-line progress when the live benchmark-relevant datum is unchanged.
\end{proof}

\begin{corollary}[What remains open after this theorem]
\label{cor:next}
After Theorem \ref{thm:main}, the remaining honest burden is no longer to justify the exact-shell-transport-section-generating substrate sub-section dynamics themselves. Those dynamics are now the canonical effective output of a sharper micro-section class. What remains open is why the sub-section-generating substrate micro-section dynamics should hold, or whether an even sharper theorem can remove that still deeper micro-section-level freedom while preserving the same benchmark one-metric package and claim boundary.
\end{corollary}

\begin{proof}
Theorem \ref{thm:main} removes the current sub-section package as the unexplained stopping point on the active line. The next honest burden therefore lies strictly below it, unless a smaller theorem-level datum inside the micro-section package can first be isolated and proved sufficient.
\end{proof}

\section{Conclusion}

The positive result of this note is narrower than a completed first-principles derivation and stronger than another same-output relay of the current section core. The current active-primary theorem had already shown that the observer-relevant sub-pre-core exact-shell transport-section core is the canonical output of a deeper exact-shell-transport-section-generating substrate sub-section dynamics package. The present theorem then removes the next unexplained stopping point by showing that this current sub-section package itself is the effective projected output of a more primitive sub-section-generating substrate micro-section dynamics.

The scope remains disciplined. The benchmark package is unchanged: the one operative metric is still exactly \(\CompletedMetric\), the ordinary matter route is unchanged, there is no second operative metric, and the low-energy matter law is not enlarged. Nothing here upgrades adoption to completed derivation or licenses stronger promotion language.

The open burden therefore moves one honest step further again. The next benchmark-facing manuscript should derive or justify the micro-section dynamics from still more primitive substrate structure, unless the sharper honest gain available instead is a genuinely smaller theorem-level collapse inside the observer-relevant micro-section package. Until then, the present result should be read for what it is: a deeper sub-section-scope derivation on the active line of \TheoryName{}, not a claim that the full first-principles program is finished.

\end{document}
